%% =============================================================================
%% CPS slot for the final -- replaces Binary Decision Trees (q/bdt.tex, 30 pts,
%% 509 lines), which was too difficult and whose `compress` reference solution
%% did not typecheck (the `val partition_bstrs = partition (...)` binding hit the
%% value restriction).
%%
%% Distinct from the practice exam's CPS question (draft_fill.tex), which is
%% search WITH BACKTRACKING over a tree. This one is about the other thing CPS
%% buys you: DISCARDING the continuation to abandon pending work.
%%
%% All three parts verified in SML/NJ:
%%   sumTree t1 = 9   prodTree t1 = 24   prodTree t2 (has a 0) = 0
%%   prodTree' t1 = 24   prodTree' t2 = 0
%% and instrumenting the multiply shows ZERO multiplications occur when a 0 is
%% found, i.e. the short-circuit is real and not merely arithmetic.
%% =============================================================================

\titledquestion{Short Circuits}

\textbf{Short Circuits}

Throughout this problem we work with the type of integer trees:
\begin{codeblock}
  datatype tree = Empty | Node of tree * int * tree
\end{codeblock}

\begin{parts}

  \part[6]\hfill

    As a warm-up, write \code{sumTree} in continuation-passing style.

    \spec
      {sumTree}
      {\code{tree -> (int -> 'a) -> 'a}}
      {\code{k} is total}
      {\code{sumTree T k} $\eeq$ \code{k s}, where \code{s} is the sum of all
      the integers in \code{T}}

    \begin{solutionorbox}[14em]
      \begin{codeblock}
        fun sumTree Empty k = k 0
          | sumTree (Node (l, x, r)) k =
              sumTree l (fn sl =>
                sumTree r (fn sr =>
                  k (sl + x + sr)))
      \end{codeblock}

      The left subtree is summed first; its continuation receives that sum and
      sums the right subtree; and \textit{that} continuation has both sums in
      hand, so it can add in \code{x} and hand the total to \code{k}.
    \end{solutionorbox}

  \newpage

  \part[10]\hfill

    Now the interesting one. Write \code{prodTree}, which multiplies the
    integers in a tree --- but which stops as soon as it meets a \code{0},
    performing \textbf{no further multiplications}.

    \spec
      {prodTree}
      {\code{tree -> (int -> int) -> int}}
      {\code{k} is total}
      {\code{prodTree T k} $\eeq$ \code{k p} if no node of \code{T} holds
      \code{0}, where \code{p} is the product of the integers in \code{T}; and
      \code{prodTree T k} $\eeq$ \code{0} otherwise}

    \begin{constraint}
      Your solution must not multiply anything once a \code{0} has been found.
      In particular, do not compute the product and then check whether it is
      \code{0}.
    \end{constraint}

    \begin{solutionorbox}[18em]
      \begin{codeblock}
        fun prodTree Empty k = k 1
          | prodTree (Node (l, x, r)) k =
              if x = 0 then 0
              else
                prodTree l (fn pl =>
                  prodTree r (fn pr =>
                    k (pl * x * pr)))
      \end{codeblock}

      The recursive case is the same shape as \code{sumTree}. The new idea is the
      \code{0} case: instead of calling \code{k} on some value, we
      \textbf{discard \code{k} entirely} and return \code{0} directly.

      That is what makes the short circuit work. Everything still to be done ---
      every pending multiplication in every enclosing subtree --- lives inside
      \code{k}, so throwing \code{k} away throws all of it away at once. A
      direct-style version cannot do this: its pending work sits on the call
      stack, where the function has no way to reach it, so each recursive call
      must return and be multiplied on the way out.
    \end{solutionorbox}

  \newpage

  \part[6]\hfill

    Notice that \code{prodTree}'s type forced the answer type to be \code{int},
    since the \code{0} case returns an \code{int} directly. Generalize it by
    taking a \textbf{second} continuation to be used when a \code{0} is found.

    \spec
      {prodTree'}
      {\code{tree -> (int -> 'a) -> (unit -> 'a) -> 'a}}
      {\code{sc} and \code{fc} are total}
      {\code{prodTree' T sc fc} $\eeq$ \code{sc p} if no node of \code{T} holds
      \code{0}, where \code{p} is the product of the integers in \code{T}; and
      \code{prodTree' T sc fc} $\eeq$ \code{fc ()} otherwise}

    Then complete the following, so that it behaves exactly like \code{prodTree}:

    \begin{codeblock}
      fun prodTree T k = prodTree' T
    \end{codeblock}
    \vspace{-0.4cm}
    \begin{solutionorbox}[6em]
      \begin{codeblock}
        k (fn () => 0)
      \end{codeblock}
    \end{solutionorbox}

    \begin{solutionorbox}[16em]
      \begin{codeblock}
        fun prodTree' Empty sc fc = sc 1
          | prodTree' (Node (l, x, r)) sc fc =
              if x = 0 then fc ()
              else
                prodTree' l (fn pl =>
                  prodTree' r (fn pr =>
                    sc (pl * x * pr)) fc) fc
      \end{codeblock}

      \code{fc} replaces the hard-coded \code{0}, which is what frees the answer
      type to be any \code{'a}. Note that \code{fc} must be threaded through
      \textbf{both} recursive calls --- a \code{0} anywhere in either subtree has
      to reach the same handler.

      Recovering \code{prodTree} is then just a matter of supplying the two
      continuations it had implicitly: \code{k} for success, and
      \code{fn () => 0} for the \code{0} case.
    \end{solutionorbox}

\end{parts}
